\setcounter{section}{4}


  \zwischenueberschrift{Übungsaufgaben}


\inputaufgabe


{}


{


$\,$


\bild{ \begin{center}


\includegraphics[width=5.5cm]{ \bildeinlesungpng {Non_cohen_macaulay_scheme_thumb} {png} }


\end{center}


\bildtext {} }


\bildlizenz { Non cohen macaulay scheme thumb.png } {} {Jakob.scholbach} {en.wikipedia} {CC-by-sa 3.0} {}


Finde ein
\definitionsverweis {Ideal}{}{,}
dessen
\definitionsverweis {Nullstellenmenge}{}{}
das folgende Gebilde ist.


}


{} {}


\inputaufgabe


{}


{


Es sei


\mavergleichskette


{\vergleichskette


{ V
}


{ \subseteq }{ { {\mathbb A}_{  K }^{  n   } }
}


{  }{
}


{    }{
}


{   }{
}


}
{}{}{}
eine Teilmenge, die aus endlich vielen Punkten bestehe. Zeige: $V$ ist genau dann
\definitionsverweis {irreduzibel}{}{,}
wenn $V$ einpunktig ist.


}


{} {}


\inputaufgabe


{}


{


Skizziere ein Beispiel einer
\definitionsverweis {zusammenhängenden}{}{,}
aber nicht
\definitionsverweis {irreduziblen}{}{}
\definitionsverweis {affin-algebraischen}{}{}
Teilmenge.


}


{} {}


\inputaufgabe


{}


{


$\,$


\bild{ \begin{center}


\includegraphics[width=5.5cm]{ \bildeinlesungsvg {Rectangular_hyperbola} {svg} }


\end{center}


\bildtext {} }


\bildlizenz { Rectangular hyperbola.svg } {} {Qef} {Commons} {PD} {}


Bestimme die irreduziblen Komponenten der reellen Hyperbel.


}


{} {}


\inputaufgabe


{}


{


Es sei $K$ ein
\definitionsverweis {Körper}{}{} der \definitionsverweis {Charakteristik}{}{}
$\neq 2$ und


\mavergleichskette


{\vergleichskette


{ a
}


{ \in }{ K
}


{  }{
}


{    }{
}


{   }{
}


}
{}{}{}
von $0$ verschieden. Zeige, dass das
\definitionsverweis {Polynom}{}{}


\mavergleichskettedisp


{\vergleichskette


{ X^2+Y^2+a
}


{ \in} { K[X,Y]
}


{ } {
}


{ } {
}


{ } {
}


}
{}{}{}
\definitionsverweis {irreduzibel}{}{}
ist.


}


{} {}


\inputaufgabe


{}


{


Es sei $K$ ein
\definitionsverweis {Körper}{}{}
und


\mavergleichskette


{\vergleichskette


{ {\mathfrak p}
}


{ = }{ (p)
}


{ \subseteq }{ K[X]
}


{    }{
}


{   }{
}


}
{}{}{}
ein
\definitionsverweis {Primideal}{}{.}
Zeige, dass die Verschwindungsmenge


\mavergleichskette


{\vergleichskette


{ V( {\mathfrak p} )
}


{ \subseteq }{ {\mathbb A}^{1}_{ K }
}


{  }{
}


{    }{
}


{   }{
}


}
{}{}{}
der gesamte


\mathl{{\mathbb A}^{1}_{K}}{}
\zusatzklammer {wobei irreduzibel vom Körper abhängt} {} {}
oder aber einpunktig
\zusatzklammer {und damit irreduzibel} {} {}
ist.


}


{} {}


\inputaufgabe


{}


{


Es sei $K$ ein
\definitionsverweis {endlicher Körper}{}{}
und


\mathl{{\mathfrak p}  \subset K[X_1 , \ldots , X_n]}{} ein
\definitionsverweis {Primideal}{}{.}
Zeige, dass die
\definitionsverweis {Verschwindungsmenge}{}{}


\mathl{V( {\mathfrak p} ) \subseteq { {\mathbb A}_{  K  }^{  n   } }}{} nur dann irreduzibel ist, wenn sie einpunktig ist.


}


{} {}


\inputaufgabe


{}


{


Zeige, dass das reelle Polynom


\mavergleichskettedisp


{\vergleichskette


{ P
}


{ =} { X^2(X-1)^2+Y^2 \left(  X^ 2+ (X-1)^2  \right)
}


{ \in} { \R[X,Y]
}


{ } {
}


{ } {
}


}
{}{}{}
ein
\definitionsverweis {Primpolynom}{}{}
ist, und dass die Nullstellenmenge


\mavergleichskettedisp


{\vergleichskette


{ V(P)
}


{ \subseteq} { \R^ 2
}


{ } {
}


{ } {
}


{ } {
}


}
{}{}{}
nicht leer, aber reduzibel ist.


}


{} {}


\inputaufgabe


{}


{


Berechne in ${ {\mathbb A}_{  \R  }^{  3   } }$ den Schnitt des Zylinders


\mathl{V \left(  x^2+y^2-1  \right)}{} mit der Kugel mit Mittelpunkt


\mavergleichskette


{\vergleichskette


{ P
}


{ = }{ (0,0,0)
}


{  }{
}


{    }{
}


{   }{
}


}
{}{}{}
und Radius $r$ in Abhängigkeit von $r$. Wann ist der Durchschnitt leer, wann irreduzibel?


}


{} {Man darf verwenden, dass der reelle Kreis irreduzibel ist.}


Die folgende Aufgabe bedeutet, dass der Durchschnitt von zwei Graphen zu Polynomen in einer Variablen mit dem Durchschnitt ihrer Differenz mit der $x$-Achse nicht nur punktweise, sondern auch algebraisch übereinstimmt. Siehe auch

Aufgabe 26.3
.


\inputaufgabegibtloesung


{}


{


Es seien


\mavergleichskette


{\vergleichskette


{ F,G
}


{ \in }{ K[X]
}


{  }{
}


{    }{
}


{   }{
}


}
{}{}{}
\definitionsverweis {Polynome}{}{}
in einer Variablen über einem
\definitionsverweis {Körper}{}{}
$K$. Zeige, dass eine
$K$-\definitionsverweis {Algebraisomorphie}{}{}


\mavergleichskettedisp


{\vergleichskette


{ K[X,Y]/(Y-F,Y-G)
}


{ \cong} { K[X]/(F-G)
}


{ } {
}


{ } {
}


{ } {
}


}
{}{}{}
vorliegt.


}


{} {}


\inputaufgabegibtloesung


{}


{


Es seien zwei verschiedene Kreise in der Ebene gegeben, die durch die Kreisgleichungen
\mathkor {} {F} {bzw.} {G} {}
gegeben seien.
\aufzaehlungzweiabc{Zeige, dass der
\definitionsverweis {Restklassenring}{}{}


\mathl{K[X,Y]/(F,G)}{}
\definitionsverweis {isomorph}{}{}
zu


\mathl{K[X,Y]/(F,H)}{} ist, wobei $H$ den Grad $\leq 1$ besitzt.
}{Zeige, dass der Restklassenring


\mathl{K[X,Y]/(F,G)}{} isomorph zu einem Ring der Form


\mathl{K[U]/(Q)}{} ist mit einem Polynom


\mavergleichskette


{\vergleichskette


{ Q
}


{ \in }{ K[U]
}


{  }{
}


{    }{
}


{   }{
}


}
{}{}{}
vom Grad $\leq 2$.
}


}


{} {}


\inputaufgabegibtloesung


{}


{


Es sei $R$ ein
\definitionsverweis {kommutativer Ring}{}{}
und


\mathl{R[X]}{} der
\definitionsverweis {Polynomring}{}{}
über $R$. Es sei


\mavergleichskette


{\vergleichskette


{ {\mathfrak a}
}


{ \subseteq }{ R[X]
}


{  }{
}


{    }{
}


{   }{
}


}
{}{}{}
ein
\definitionsverweis {Ideal}{}{}
mit
\definitionsverweis {Erzeugern}{}{}


\mavergleichskettedisp


{\vergleichskette


{ {\mathfrak a}
}


{ =} { \left(  F_0, F_1 , \ldots , F_n  \right)
}


{ } {
}


{ } {
}


{ } {
}


}
{}{}{,}
wobei


\mavergleichskette


{\vergleichskette


{F_0
}


{ = }{ X-r
}


{  }{
}


{    }{
}


{   }{
}


}
{}{}{}
mit


\mavergleichskette


{\vergleichskette


{ r
}


{ \in }{ R
}


{  }{
}


{    }{
}


{   }{
}


}
{}{}{}
sei. Für


\mavergleichskette


{\vergleichskette


{ i
}


{ \geq }{ 1
}


{  }{
}


{    }{
}


{   }{
}


}
{}{}{}
seien


\mathl{G_i}{} die Elemente aus $R$, die entstehen, wenn man in $F_i$ die Variable $X$ durch $r$ ersetzt. Zeige, dass eine
\definitionsverweis {Ringisomorphie}{}{}
der
\definitionsverweis {Restklassenringe}{}{}


\mavergleichskettedisp


{\vergleichskette


{ R[X] / {\mathfrak a}
}


{ \cong} { R/ \left(  G_1 , \ldots , G_n  \right)
}


{ } {
}


{ } {
}


{ } {
}


}
{}{}{}
vorliegt.


}


{} {}


\inputaufgabe


{}


{


Betrachte die Menge der reellen Zahlen $\R$ mit der metrischen Topologie. Ist $\R$
\definitionsverweis {irreduzibel}{}{?}


}


{} {}


\inputaufgabegibtloesung


{}


{


Es sei $p$ eine
\definitionsverweis {Primzahl}{}{}
und $\Z/( p )$ der zugehörige
\definitionsverweis {Restklassenkörper}{}{.}
Zeige: Jede Quadrik der Form


\mavergleichskettedisp


{\vergleichskette


{ F
}


{ =} { aX^2+bY^2+c
}


{ =} { 0
}


{ } {
}


{ } {
}


}
{}{}{}
mit


\mavergleichskette


{\vergleichskette


{ a,b
}


{ \neq }{ 0
}


{  }{
}


{    }{
}


{   }{
}


}
{}{}{}
hat mindestens eine Lösung in $\Z/( p )$.


}


{} {}


\inputaufgabegibtloesung


{}


{


Es sei ${\mathfrak a}$ ein
\definitionsverweis {Ideal}{}{}
in einem
\definitionsverweis {kommutativen Ring}{}{}
$R$. Zeige, dass ${\mathfrak a}$ genau dann ein
\definitionsverweis {Primideal}{}{}
ist, wenn ${\mathfrak a}$ der
\definitionsverweis {Kern}{}{}
eines
\definitionsverweis {Ringhomomorphismus}{}{}
\maabb {\varphi} { R } { K
} {}
in einen
\definitionsverweis {Körper}{}{}
$K$ ist.


}


{} {}


\inputaufgabe


{}


{


Zeige, dass ein
\definitionsverweis {maximales Ideal}{}{}
${\mathfrak m}$ in einem
\definitionsverweis {kommutativen Ring}{}{}
$R$ ein
\definitionsverweis {Primideal}{}{}
ist.


}


{} {}


\inputaufgabegibtloesung


{}


{


Es sei $R$ ein
\definitionsverweis {kommutativer Ring}{}{}
und ${\mathfrak p}$ ein
\definitionsverweis {Ideal}{}{.}
Zeige, dass ${\mathfrak p}$ genau dann ein
\definitionsverweis {Primideal}{}{}
ist, wenn der
\definitionsverweis {Restklassenring}{}{}
$R/{\mathfrak p}$ ein
\definitionsverweis {Integritätsbereich}{}{}
ist.


}


{} {}


\inputaufgabe


{}


{


Es sei


\mavergleichskette


{\vergleichskette


{ {\mathfrak p}
}


{ \subseteq }{ R
}


{  }{
}


{    }{
}


{   }{
}


}
{}{}{}
ein
\definitionsverweis {Primideal}{}{}
in einem
\definitionsverweis {kommutativen Ring}{}{}
$R$. Zeige, dass aus einer Inklusionsbeziehung


\mavergleichskettedisp


{\vergleichskette


{ {\mathfrak a} \cap {\mathfrak b}
}


{ \subseteq} { {\mathfrak p}
}


{ } {
}


{ } {
}


{ } {
}


}
{}{}{}
die Inklusion


\mavergleichskette


{\vergleichskette


{ {\mathfrak a}
}


{ \subseteq }{ {\mathfrak p}
}


{  }{
}


{    }{
}


{   }{
}


}
{}{}{}
oder


\mavergleichskette


{\vergleichskette


{ {\mathfrak b}
}


{ \subseteq }{ {\mathfrak p}
}


{  }{
}


{    }{
}


{   }{
}


}
{}{}{}
folgt.


}


{} {}


\inputaufgabe


{}


{


Es seien $R$ und $S$
\definitionsverweis {kommutative Ringe}{}{}
und sei
\maabb {\varphi} { R } { S
} {}
ein
\definitionsverweis {Ringhomomorphismus}{}{.}
Es sei ${\mathfrak p}$ ein
\definitionsverweis {Primideal}{}{}
in $S$. Zeige, dass das Urbild


\mathl{\varphi^{-1}( {\mathfrak p} )}{} ein Primideal in $R$ ist.


Zeige durch ein Beispiel, dass das Urbild eines
\definitionsverweis {maximalen Ideales}{}{}
kein maximales Ideal sein muss.


}


{} {}


\inputaufgabe


{}


{


Es sei $K$ ein
\definitionsverweis {Körper}{}{} und


\mavergleichskette


{\vergleichskette


{ L
}


{ = }{ K(X)
}


{  }{
}


{    }{
}


{   }{
}


}
{}{}{}
der
\definitionsverweis {Quotientenkörper}{}{}
des
\definitionsverweis {Polynomrings}{}{}
$K[X]$. Zeige, dass es unendlich viele Zwischenkörper zwischen
\mathkor {} {K} {und} {L} {}
gibt.


}


{} {}


\inputaufgabe


{}


{


Erkläre, wo der
Beweis
 zu

Satz 4.8

zusammenbricht, wenn man ihn auf mehr als zwei Variablen ausdehnen will.


}


{} {}


\inputaufgabe


{}


{


Es seien $P(X)$ und $Q(Y)$ nichtkonstante Polynome in der einen angegebenen Variablen. Man gebe eine Abschätzung
\zusatzklammer {unter welcher Bedingung} {?} {}
für die Anzahl der Schnittpunkte der beiden Kurven


\mathl{V(Y-P(X))}{} und


\mathl{V(X-Q(Y))}{.}


}


{} {}


In den folgenden Aufgaben werden die Begriffe \stichwort {abgeschlossene Abbildung} {} und \stichwort {offene Abbildung} {} verwendet.


Eine
\definitionsverweis {stetige Abbildung}{}{}
\maabbdisp {f} { X } { Y
} {}
zwischen
\definitionsverweis {topologischen Räumen}{}{}
$X$ und $Y$ heißt \definitionswort {abgeschlossen}{,} wenn Bilder von
\definitionsverweis {abgeschlossenen Mengen}{}{}
wieder abgeschlossen sind.


Sie heißt \stichwort {offen} {,} wenn Bilder von offenen Mengen wieder offen sind.


\inputaufgabe


{}


{


Zeige, dass die Projektion
\maabbeledisp {} { {\mathbb A}^{2}_{ K } } { {\mathbb A}^{1}_{K}
} { (x,y) } { x
} {,}
nicht
\definitionsverweis {abgeschlossen}{}{}
in der \definitionsverweis {Zariski-Topologie}{}{}
ist.


}


{} {}


\inputaufgabe


{}


{


Zeige, dass eine ebene algebraische Kurve über den komplexen Zahlen $\mathbb C$ nicht \definitionsverweis {kompakt}{}{} in der metrischen Topologie ist.


}


{} {}


  \zwischenueberschrift{Aufgaben zum Abgeben}


\inputaufgabe


{6}


{


Es sei $p$ eine
\definitionsverweis {Primzahl}{}{}
$\geq 3$ und


\mavergleichskette


{\vergleichskette


{ K
}


{ = }{ \Z/( p )
}


{  }{
}


{    }{
}


{   }{
}


}
{}{}{}
der zugehörige
\definitionsverweis {Restklassenkörper}{}{.}
Es sei ein Polynom


\mavergleichskette


{\vergleichskette


{ F
}


{ \in }{ \Z/( p ) [X,Y]
}


{  }{
}


{    }{
}


{   }{
}


}
{}{}{}
der Form


\mavergleichskettedisp


{\vergleichskette


{ F
}


{ =} { \alpha X^2+\beta XY + \gamma Y^2 + \delta X + \epsilon Y + \eta
}


{ } {
}


{ } {
}


{ } {
}


}
{}{}{}
gegeben. Zeige, dass für das zugehörige Nullstellengebilde


\mavergleichskette


{\vergleichskette


{ V(F)
}


{ \subseteq }{ {\mathbb A}^{2}_{ K }
}


{  }{
}


{    }{
}


{   }{
}


}
{}{}{}

(wenn $\alpha, \beta, \gamma$ nicht alle $0$ sind, so ist das eine Quadrik) die folgenden drei Alternativen bestehen.
\aufzaehlungdrei{ $V(F)$ besitzt mindestens einen Punkt.
}{


\mavergleichskette


{\vergleichskette


{ F
}


{ = }{ c
}


{  }{
}


{    }{
}


{   }{
}


}
{}{}{}
mit einer Konstanten


\mavergleichskette


{\vergleichskette


{ c
}


{ \neq }{ 0
}


{  }{
}


{    }{
}


{   }{
}


}
{}{}{.}
}{Es gibt eine Variablentransformation derart, dass das Polynom in den neuen Koordinaten die Gestalt $Z^2-u$ mit einem Nichtquadrat


\mavergleichskette


{\vergleichskette


{ u
}


{ \in }{ \Z/( p )
}


{  }{
}


{    }{
}


{   }{
}


}
{}{}{}
besitzt.
}


}


{} {Tipp: Für einen wichtigen Fall ist

Aufgabe 9.20 (Zahlentheorie (Osnabrück 2025))

hilfreich.}


\inputaufgabe


{3}


{


Es sei $V$ eine irreduzible, affin-algebraische Menge mit mindestens zwei Punkten und seien


\mavergleichskette


{\vergleichskette


{ P_1 , \ldots ,  P_m
}


{ \in }{ V
}


{  }{
}


{    }{
}


{   }{
}


}
{}{}{}
endlich viele Punkte darin. Zeige, dass dann auch


\mathl{V \setminus \{ P_1 , \ldots , P_m \}}{}
\zusatzklammer {in der induzierten Topologie} {} {}
irreduzibel ist.


}


{} {}


\inputaufgabe


{4}


{


Es sei $R$ ein
\definitionsverweis {faktorieller Integritätsbereich}{}{}
mit
\definitionsverweis {Quotientenkörper}{}{}
$Q(R)$. Zeige: Wenn


\mavergleichskette


{\vergleichskette


{ F,G
}


{ \in }{ R[X]
}


{  }{
}


{    }{
}


{   }{
}


}
{}{}{}
keinen gemeinsamen
\zusatzklammer {nichtkonstanten} {} {}
Teiler besitzen, so besitzen sie aufgefasst in


\mathl{Q(R)[X]}{} ebenfalls keinen gemeinsamen Teiler.


}


{\zusatzklammer {Man darf sich auf Hauptidealbereiche $R$  beschränken.} {} {}} {}


\inputaufgabe


{3}


{


Es sei


\mavergleichskette


{\vergleichskette


{ K
}


{ = }{ \Q
}


{  }{
}


{    }{
}


{   }{
}


}
{}{}{}
der
\definitionsverweis {Körper}{}{}
der
\definitionsverweis {rationalen Zahlen}{}{.}
Begründe, ob


\mavergleichskettedisp


{\vergleichskette


{ V \left(  X^2+Y^2-1  \right)
}


{ \subseteq} { {\mathbb A}^{2}_{\Q}
}


{ } {
}


{ } {
}


{ } {
}


}
{}{}{}
\definitionsverweis {irreduzibel}{}{}
ist oder nicht.


}


{} {Tipp: Man verwende

Aufgabe 1.28

und

Korollar 4.9
.}


\inputaufgabe


{4}


{


Zeige, dass die Projektion
\maabbeledisp {} {  {\mathbb A}^{2}_{K} } { {\mathbb A}^{1}_{K}
} {(x,y)} {x
} {,}
\definitionsverweis {offen}{}{}
in der \definitionsverweis {Zariski-Topologie}{}{}
ist.


}


{} {}


\inputaufgabe


{4}


{


Zeige, dass die affine Ebene $\mathbb A^2_K$ mit der
\definitionsverweis {Zariski-Topologie}{}{}
\definitionsverweis {kompakt}{}{}
ist.


}


{} {}
